\subsection{Communication as Action}
\label{sec:communication-as-action}
\slidesrc{Pragmatics: 4--10}

These slides come from another course (on agents) and give a very quick introduction to
``communication as an action''.

\subsubsection{From true/false sentences to actions}
\slidesrc{Pragmatics: 5--8}

\begin{nbdef}{Communication}
\textbf{Communication}, in general, is the \emph{intentional} exchange of information
brought about by the production and perception of \emph{signs} drawn from a
\emph{shared system of conventional signs}.
\end{nbdef}

Two views of language:
\begin{itemize}[nosep]
  \item the \textbf{``classical'' view} (before 1953): language consists of sentences
        that are \textbf{true or false} (as in logic; this is the view behind
        model-theoretic semantics, Section~\ref{sec:semantic-parsing});
  \item the \textbf{``modern'' view} (after 1953): \textbf{language is a form of
        action}.
\end{itemize}
The turning point is marked by three philosophers.

\paragraph{Wittgenstein (1953), \emph{Philosophical Investigations}.}
According to Wittgenstein:
\begin{itemize}[nosep]
  \item the structure of reality determines the structure of language;
  \item \textbf{meaning is determined by use}, in the context of a
        \textbf{``language-game''};
  \item words are not defined by reference to the objects or things they designate in
        the external world, nor by the thoughts, ideas or mental representations one
        might associate with them, but by \textbf{how they are used} in effective,
        ordinary communication.
\end{itemize}

\begin{nbintuition}[language-games]
Think of the word \emph{``Slab!''} shouted by a builder to an assistant: it does not
describe a slab, it means ``bring me a slab'', and it means that only inside the
``game'' of building together. Every use of language is a move in some game with its own
rules: giving orders, asking, telling a joke, praying. This is the same Wittgenstein of
``meaning is use'' at the root of distributional semantics
(Section~\ref{sec:distributional-hypothesis}).
\end{nbintuition}

\begin{nbwarning}[two Wittgensteins]
The first bullet (``the structure of reality determines the structure of language'') is
the idea of the \emph{early} Wittgenstein (\emph{Tractatus}, 1921), in which a sentence is
a picture of a fact. The \emph{Philosophical Investigations} (1953) largely abandon it in
favour of meaning as use. The slide lists both; remember that the 1953 work is the one of
language-games.
\end{nbwarning}

\paragraph{Austin (1962), \emph{How to Do Things with Words}.}
\textbf{In saying something, we do something}: when a minister joins two people in
marriage saying ``I now pronounce you husband and wife'', the sentence does not describe
a marriage, it \emph{makes} it (Section~\ref{sec:course-map}). Such sentences are
\textbf{performatives}.

\paragraph{Searle (1969), \emph{Speech Act Theory}.}
\begin{itemize}[nosep]
  \item Each \textbf{illocution} (the act performed \emph{in} saying something) can be
        described in terms of what it is attempting to do. An \emph{assertion} counts as
        an undertaking that something really is the case; a \emph{question} counts as an
        attempt to elicit some information; \emph{thanking} someone counts as an
        expression of gratitude.
  \item What matters is the assumed \textbf{intent of the speaker}, the
        \textbf{intentionality} of the sentence.
\end{itemize}

\notinslides
\begin{nbdef}{Speech act: three levels (Austin)}
\begin{itemize}[nosep]
  \item \textbf{Locution}: the act of saying something, the words with their literal
        meaning (``It's cold in here'').
  \item \textbf{Illocution}: the act performed \emph{in} saying it, the speaker's
        intention (a request to close the window).
  \item \textbf{Perlocution}: the effect produced \emph{by} saying it on the hearer
        (the hearer gets up and closes the window).
\end{itemize}
\end{nbdef}

\subsubsection{Speech acts and the stages of communication}
\slidesrc{Pragmatics: 9--10}

\begin{center}
\begin{tikzpicture}[b/.style={font=\small}]
  \node[draw=nbred, rounded corners=6pt, minimum width=6.2cm, minimum height=1.1cm] (sit) {};
  \node[font=\tiny, anchor=north west, text=nbred] at (sit.north west) {SITUATION};
  \node[b, text=nbblue] (s) at (-2.0,-0.1) {Speaker};
  \node[b, text=nbred] (u) at (0.1,-0.1) {Utterance};
  \node[b, text=nbblue] (h) at (2.2,-0.1) {Hearer};
  \draw[-{Stealth}] (s) -- (u); \draw[-{Stealth}] (u) -- (h);
\end{tikzpicture}
\captionof{figure}{A speech act: a speaker produces an utterance for a hearer, inside a
situation.}
\end{center}

\textbf{Speech acts achieve the speaker's goals}. The examples come from the Wumpus world,
the cave game used for the grammar of Section~\ref{sec:grammars-parsing}:
\begin{center}
\small
\begin{tabular}{@{}l l@{}}
\toprule
\textbf{act} & \textbf{example} \\
\midrule
Inform & ``There's a pit in front of you'' \\
Query & ``Can you see the gold?'' \\
Command & ``Pick it up'' \\
Promise & ``I'll share the gold with you'' \\
Acknowledge & ``OK'' \\
\bottomrule
\end{tabular}
\end{center}
\textbf{Speech act planning} (deciding what to say to reach a goal) requires knowledge
of the \emph{situation}, of the \emph{semantic and syntactic conventions}, and of the
\emph{hearer's goals, knowledge base and rationality}.

\begin{center}
\small
\begin{tabular}{@{}l l l@{}}
\toprule
 & \textbf{stage} & \textbf{what happens} \\
\midrule
\multirow{3}{*}{\textcolor{nbblue}{speaker S}}
  & Intention & S wants to inform H that $P$ \\
  & Generation & S selects words $W$ to express $P$ in context $C$ \\
  & Synthesis & S utters words $W$ \\
\midrule
\multirow{4}{*}{\textcolor{nbblue}{hearer H}}
  & Perception & H perceives $W'$ in context $C'$ \\
  & Analysis & H infers possible meanings $P_1, \ldots, P_n$ \\
  & Disambiguation & H infers the intended meaning $P_i$ \\
  & Incorporation & H incorporates $P_i$ into its knowledge base \\
\bottomrule
\end{tabular}
\captionof{table}{The stages in communication.}
\label{tab:communication-stages}
\end{center}

\begin{nbintuition}[where things can go wrong]
Notice the primes: the hearer perceives $W'$ in context $C'$, not necessarily the same
$W$ and $C$ of the speaker (noise, mishearing, a different idea of the situation). The
analysis gives \emph{several} possible meanings, and disambiguation picks one, maybe the
wrong one. Every stage maps onto a part of NLP: generation and synthesis are natural
language generation and speech synthesis; perception is speech recognition or reading;
analysis is syntax and semantics (parsing, Sections~\ref{sec:sentence-level-syntax}
to~\ref{sec:lexical-semantics}); disambiguation is WSD and intent recognition;
incorporation is updating a knowledge base, e.g.\ an ontology
(Section~\ref{sec:ontologies}).
\end{nbintuition}
